\section{Глава~\ref{sec:serocomputer}. \nt{СЕРО}-нейроны}

Свойства самих \nt{СЕРО}-нейронов (ритм, автоторможение, знак по рецептору, подсистемы) измерены прямо; вычисления главы --- регулятор, фильтр, логика --- следуют из её уравнений; $\tau_c$, коэффициент диффузии серотонина и число мест выброса --- оценки; роль петли как регулятора уровня в мозге --- гипотеза (в ядре шва автоторможение точечное и даёт лишь короткие паузы), а роль \nt{СЕРО} как горизонта --- гипотеза с поведенческими опытами в её пользу.

{\footnotesize
\setlength{\tabcolsep}{3pt}\renewcommand{\arraystretch}{1.15}
\begin{longtable}{>{\raggedright}p{0.5cm}>{\raggedright}p{4.0cm}>{\raggedright}p{5.6cm}>{\raggedright}p{2.3cm}>{\raggedright\arraybackslash}p{1.7cm}}
\toprule
№ & Утверждение & Опыт и что измерено & Источник & Статус \\
\midrule\endfirsthead
\toprule
№ & Утверждение & Опыт и что измерено & Источник & Статус \\
\midrule\endhead
1 & Знак действия серотонина выбирает рецептор получателя (утверждение~\ref{prop:receptor}) & Рецепторы серотонина разделены на семейства по фармакологии и механизму: 5-HT$_{1A}$ через G-белок открывает калиевые каналы и тормозит клетку, 5-HT$_{2A}$ и ионотропный 5-HT$_3$ возбуждают. Один медиатор даёт противоположные ответы в разных клетках. & \cite{BarnesSharp1999} & измерено \\
2 & \nt{СЕРО}-нейрон в бодрствовании ровно выдаёт несколько импульсов в секунду & Запись нейронов дорсального ядра шва у свободно движущихся кошек: медленный ритмичный разряд $2.8\pm0.2$ имп/с в спокойном бодрствовании; в медленном сне частота падает на $34$--$68\,\%$, в быстром --- на $98\,\%$. Под наркозом у крыс --- $1.7\pm0.5$ имп/с, очень регулярно. & \cite{TrulsonJacobs1979, GagnonParent2014, JacobsAzmitia1992} & измерено \\
3 & \nt{СЕРО}-нейрон --- ритмоводитель: толчок на каждый импульс не нужен, но нужен ровный фоновый приток (в срезе --- норадреналин через $\alpha_1$-рецепторы, в организме --- от \nt{НОРА}) & В срезе мозга клетки разряжаются медленно и ровно, после каждого импульса --- следовая гиперполяризация на $200$--$800\ms$. Спонтанно работающих клеток в срезе меньше, чем в организме: ровному ритму нужен тонический вход норадреналина через $\alpha_1$-рецепторы, блокада которых его гасит. & \cite{VandermaelenAghajanian1983} & измерено \\
4 & Почти все рецепторы метаботропные; ответ медленный: нарастает и спадает в пределах примерно секунды & Все семейства рецепторов, кроме 5-HT$_3$, связаны с G-белками. В срезе вызванный выброс серотонина даёт через 5-HT$_{1A}$ тормозящий ток, который нарастает и спадает в пределах секунды. & \cite{BarnesSharp1999, CourtneyFord2016} & измерено \\
5 & В областях назначения серотонин выходит в объём, а не только в щель; в самом ядре шва торможение соседей идёт точка-точка & Быстрая циклическая вольтамперометрия в срезах: выброс на один импульс одинаков в области с синапсами и в ядре шва, где синапсов почти нет; не зависит от блокады захвата и рецепторов; молекул на окончание намного меньше, чем переносчиков и рецепторов, а внеклеточная концентрация совпадает со сродством главного рецептора. В ядре шва торможение через 5-HT$_{1A}$ идёт точка-точка: переносчик не даёт серотонину накопиться вне щели. & \cite{Bunin1998, CourtneyFord2016} & измерено \\
6 & Концентрация по~\eqref{eq:seroneuron} убывает из-за откачки; $\tau_c$ порядка секунды --- оценка & Вольтамперометрия в живом мозге: внеклеточный серотонин ограничивают прежде всего обратный захват и разрушение, а не синтез. Прямого измерения $\tau_c$ в цитируемых работах нет; порядок величины --- оценка главы. & \cite{Hashemi2012serotonin} & измерено; $\tau_c$ --- оценка \\
7 & НЕ и ИЛИ-НЕ на одном ритмоводителе; полнота \nt{СЕРО}-логики (утверждение~\ref{prop:serocomplete}, рис.~\ref{fig:serogates}) & Следует из~\eqref{eq:seroneuron}: тормозимый ритмоводитель --- ИЛИ-НЕ, а ИЛИ-НЕ функционально полон. Опыта, где из \nt{СЕРО}-нейронов собрана логика, нет; условие раздельных объёмов в мозге не выполняется. & вывод в главе & теория \\
8 & Серотонин тормозит сам \nt{СЕРО}-нейрон через рецепторы на нём & У наркотизированных крыс агонисты 5-HT$_{1A}$ внутривенно и ионофоретически тормозят разряд нейронов дорсального шва с той же зависимостью доза--ответ, что и сам серотонин; агонисты 5-HT$_{1B}$ почти не действуют. В срезе эндогенный серотонин соседних клеток даёт через 5-HT$_{1A}$ ток и короткую паузу в разряде, не меняя средней частоты. & \cite{Sprouse1987, CourtneyFord2016} & измерено \\
9 & В модели петля через общий объём --- регулятор уровня: разброс и постоянная времени уменьшаются в $1+G$ раз~\eqref{eq:feedback}; что петля так работает в мозге, --- гипотеза & Классическая формула усилителя с отрицательной обратной связью; в главе выведена заново. Усиление петли $G$ у \nt{СЕРО}-системы не измерено. Измерено: в срезе автоторможение идёт точка-точка, даёт короткую паузу и не меняет средней частоты. & \cite{Black1934feedback, CourtneyFord2016} & теория; в мозге --- гипотеза \\
10 & Концентрация --- скользящее среднее частоты, фильтр нижних частот~\eqref{eq:lowpass}, рис.~\ref{fig:lowpass} & Решение линейного уравнения~\eqref{eq:seroneuron}; рисунок --- расчёт по этой формуле при $\tau_c=1\,\text{с}$. Отдельной модели в \texttt{models/} нет. & вывод в главе & теория \\
11 & Извилистость щелей замедляет диффузию малой молекулы примерно в $2.6$ раза & Метод точечного источника: ионофорез тетраметиламмония и запись его концентрации ионоселективным электродом в срезах и в живом мозге. Извилистость $\lambda\approx1.6$, то есть эффективный $D$ меньше свободного в $\lambda^2\approx2.6$ раза; доля межклеточного объёма около $0.2$. Коэффициент диффузии самого серотонина в ткани в этих работах не измерялся; в главе он --- оценка. & \cite{Sykova2008} & измерено \\
12 & Длина независимого «провода» $\ell\sim\sqrt{D\tau}\approx20$~мкм~\eqref{eq:diffusionlength}; число проводов ограничено числом таких областей (утверждение~\ref{prop:volumewires}) & Закон диффузии и подстановка оценок $D$ и $\tau$ (строки~6 и~11); утверждение~\ref{prop:volumewires} --- следствие. Прямого опыта, считающего независимые каналы объёмной передачи, нет. & вывод в главе & теория \\
13 & Выход одного \nt{СЕРО}-нейрона раскинут по огромным участкам мозга; мест выброса по оценке $\sim10^5$ & Полная реконструкция отдельных нейронов дорсального шва у крысы: все восходящие аксоны сильно ветвятся, общая длина аксона одной клетки до $18.7$~см, ветви одной клетки доходят до стриатума, префронтальной коры и миндалины. Число мест выброса на клетку прямо не подсчитано. & \cite{GagnonParent2014} & измерено; $10^5$ --- оценка \\
14 & \nt{СЕРО}-нейроны делятся на несколько подсистем с разными выходами и ответами & Вирусно-генетическое мечение у мышей: нейроны, идущие к миндалине и к лобной коре, почти не пересекаются по ветвлению, получают разные входы и противоположно отвечают на неприятный стимул; включение и выключение каждой группы по-разному меняет поведение. & \cite{Ren2018} & измерено \\
15 & Условие терпения $D<\tau_\gamma\ln(R/r_0)$~\eqref{eq:patience} при экспоненциальном обесценивании; при измеренном, гиперболическом, --- $D<\tau_\gamma(R/r_0-1)$ & Следует из обесценивания $e^{-D/\tau_\gamma}$ или $1/(1+D/\tau_\gamma)$; вывод в главе. Обезьяны, выбор на задержках в секунды: обесценивание ближе к гиперболе, чем к экспоненте; ответ \nt{ДОФА}-нейронов на предвестник отложенной награды падает с задержкой тоже по гиперболе. & вывод в главе; \cite{KobayashiSchultz2008} & теория \\
16 & \nt{СЕРО} задаёт горизонт $\tau_\gamma$ (раздел~\ref{sec:horizon}) & Оптогенетическое включение \nt{СЕРО}-нейронов дорсального шва у мышей удлиняет ожидание отложенной награды и повышает упорство в поиске награды, ставшей редкой. У людей понижение серотонина (диета без триптофана) делает обесценивание отложенной награды круче. Как концентрация превращается в $\tau_\gamma$, неизвестно. & \cite{Doya2002, Miyazaki2014, Lottem2018, Schweighofer2008} & гипотеза \\
\bottomrule
\end{longtable}}
